\documentclass[12pt]{report}
% Template Tesi in Informatica, ispirata al template "Tesi di laurea - Universit� di Pisa" di Simone Schirinzi (https://www.overleaf.com/latex/templates/tesi-di-laurea-universita-di-pisa/rwdcqtqwftpg). 

% Carattere dimensione 12

% Per la stampa fronte-retro sostituire con:
% \documentclass[12pt, twoside]{report}

% Margini (4cm a sx, 2.5cm a dx, 2.5cm in alto, 2.5cm in basso)
\usepackage[top=2.5cm, bottom=2.5cm, left=4cm, right=2.5cm, head=21.75pt, centering]{geometry}

% Per la stampa fronte-retro sostituire con: 
% \usepackage[top=2.5cm, bottom=2.5cm, inner=4cm, outer=4cm, right=2.5cm, centering]{geometry}

% Interlinea
\linespread{1.5}

% Librerie utili
\usepackage[italian]{babel} % applicazione regole di scrittura per la lingua italiana 
\usepackage[utf8]{inputenc} % codifica UTF-8
\usepackage{scrlayer-scrpage} % stili pagina per il frontespizio
\usepackage[hidelinks]{hyperref} % collegamenti ipertestuali
\ifoot[]{}
\cfoot[]{}
\ofoot[\pagemark]{\pagemark}
\pagestyle{scrplain}
% \usepackage{mathptmx} % font Times New Roman (simile)
\usepackage{graphicx} % inserimento di immagini
\usepackage{csquotes} % per le citazioni "in blocco"
% \usepackage[backend=biber, sorting=nty, ]{biblatex} % bibliografia con pacchetto biblatex (https://ctan.org/pkg/biblatex?lang=en)
\appto{\bibsetup}{\raggedright}

\usepackage{titlesec} % per la formattazione dei titoli delle sezioni, capitoli etc.
\usepackage{float} % per il posizionamento delle immagini

\usepackage{listings} % per il codice di programmazione
% Fonte https://en.wikibooks.org/wiki/LaTeX/Source_Code_Listings. Per la lista di sintassi riconosciute.
\renewcommand{\lstlistingname}{Code}% Listing -> Codice
\usepackage{xcolor}  % stile del codice
\definecolor{mygreen}{rgb}{0,0.6,0}
\definecolor{mygray}{rgb}{0.5,0.5,0.5}
\definecolor{mymauve}{rgb}{0.58,0,0.82}
\definecolor{darkgray}{rgb}{.4,.4,.4}
\definecolor{navy}{HTML}{000080}
\definecolor{purple}{rgb}{0.65, 0.12, 0.82}
\definecolor{codepurple}{rgb}{0.58,0,0.82}
\definecolor{backcolour}{rgb}{0.95,0.95,0.92}

% Stili configurabili del codice (lslisting) 
\lstset{ %
  belowcaptionskip=0.5em,
  backgroundcolor=\color{backcolour}, % choose the background color; you must add \usepackage{color} or \usepackage{xcolor}
  basicstyle=\footnotesize, % the size of the fonts that are used for the code
  breakatwhitespace=false, % sets if automatic breaks should only happen at whitespace
  breaklines=true, % sets automatic line breaking
  captionpos=b, % sets the caption-position to bottom
  commentstyle=\color{mygreen}, % comment style
  deletekeywords={...}, % if you want to delete keywords from the given language
  escapeinside={\%*}{*)}, % if you want to add LaTeX within your code
  extendedchars=true, % lets you use non-ASCII characters; for 8-bits encodings only, does not work with UTF-8
  frame=single, % adds a frame around the code
  keepspaces=true, % keeps spaces in text, useful for keeping indentation of code (possibly needs columns=flexible)
  keywordstyle=\color{codepurple}, % keyword style
  language=Python, % the language of the code
  morekeywords={*,...}, % if you want to add more keywords to the set
  numbers=left, % where to put the line-numbers; possible values are (none, left, right)
  numbersep=5pt, % how far the line-numbers are from the code
  numberstyle=\tiny\color{mygray}, % the style that is used for the line-numbers
  rulecolor=\color{black}, % if not set, the frame-color may be changed on line-breaks within not-black text (e.g. comments (green here))
  showspaces=false, % show spaces everywhere adding particular underscores; it overrides 'showstringspaces'
  showstringspaces=false, % underline spaces within strings only
  showtabs=false, % show tabs within strings adding particular underscores
  stepnumber=1, % the step between two line-numbers. If it's 1, each line will be numbered
  stringstyle=\color{mymauve}, % string literal style
  tabsize=2, % sets default tabsize to 2 spaces
  title=\lstname % show the filename of files included with \lstinputlisting; also try caption instead of title
}

% END of listing package 

% Esempio riconoscimento sintassi JavaScript 

\lstdefinelanguage{JavaScript}{
  keywords={typeof, new, true, false, catch, function, return, null, catch, switch, var, const, let, if, in, while, do, else, case, break},
  keywordstyle=\color{blue}\bfseries,
  ndkeywords={class, export, boolean, throw, implements, import, this},
  ndkeywordstyle=\color{darkgray}\bfseries,
  identifierstyle=\color{black},
  sensitive=false,
  comment=[l]{//},
  morecomment=[s]{/*}{*/},
  commentstyle=\color{mygreen}\ttfamily,
  stringstyle=\color{red}\ttfamily,
  morestring=[b]',
  morestring=[b]"
}

\lstset{
   language=JavaScript,
   extendedchars=true,
   basicstyle=\footnotesize\ttfamily,
   showstringspaces=false,
   showspaces=false,
   numbers=left,
   numberstyle=\footnotesize,
   numbersep=9pt,
   tabsize=2,
   breaklines=true,
   showtabs=false,
   captionpos=b
}

% Formato delle intestazioni
\titleformat{\chapter}[block]
  {\normalfont\LARGE\bfseries}{\thechapter.}{0.5em}{\LARGE}
\titlespacing*{\chapter}{0pt}{-20pt}{25pt}

\begin{document}

% Frontespizio
\begin{titlepage}
  \begin{figure}
    \centering\includegraphics[scale=0.5]{immagini/cherubino_pant541.png}
  \end{figure}

  \begin{center}
    {\LARGE{ Corso di Laurea in Informatica \\ }}
    \vspace{2cm}
    {\Large { TESI DI LAUREA }}\\
    \vspace{2cm}
    {\Large { Algoritmi genetici per Explainable AI }}
  \end{center}

  \vspace{2cm}

  \begin{minipage}[t]{0.47\textwidth}
    {\large{Relatori:\\ Riccardo Guidotti}}
    \vspace{0.5cm}
    {\large{\\Correlatore:\\ Marco Danelutto}}
  \end{minipage}\hfill\begin{minipage}[t]{0.47\textwidth}\raggedleft
    {\large{Candidato: \\ Federico Bustaffa\\ }}
  \end{minipage}

  \vspace{25mm}

  \centering{\large{\bf ANNO ACCADEMICO 2024/2025 }}
\end{titlepage}
% Fine frontespizio

\tableofcontents
\thispagestyle{empty}

\listoffigures

\thispagestyle{empty}
\clearpage
\setcounter{page}{1}
\addtocontents{toc}{\protect\thispagestyle{empty}}
\addcontentsline{toc}{chapter}{Introduzione} % Capitolo non numerato

\chapter*{Introduzione}

Il lavoro di tesi verte sull'implementazione di una libreria di algoritmi
genetici, in grado di sfruttare architetture multicore per il calcolo parallelo.
L'esigenza di implementare un algoritmo con queste caratteristiche nasce dalla
necessità di ottimizzare un precedente lavoro, il quale ne fa largo uso.

Il progetto in questione è \textit{LORE}~\cite{guidotti2018LORE}, il quale si
colloca nell'ambito dell'explainable AI e che ha quindi l'obbiettivo di
\textit{spiegare} le decisioni prese da un modello di machine learning.
Attualmente \textit{LORE} sfrutta \textit{DEAP}~\cite{fortin2012DEAP}, come
libreria di algoritmi genetici, la quale offre un alto livello di modularità e
flessibilità, mettendo a disposizione dell'utente la possibilità di sfruttare
il parallelismo, dando però poco controllo sull'architettura del modello di
calcolo.

Si è quindi cercato di implementare un'architettura più ottimizzata e specifica,
in g

Le implementazioni di \textit{LORE} e \textit{DEAP}, così come la maggior parte
delle librerie di machine learning forniscono una API Python, rendendo quindi
necessario implementare una libreria con una API nello stesso linguaggio.

Per valutare qualità e correttezza sono stati svolti test su problemi classici,
come il problema del commesso viaggiatore e dello zaino o, per valutare il
comportamento su problemi a valori reali, un semplice caso di regressione
lineare. Tutti problemi in cui la qualità della soluzione generata fosse
facilmente verificabile e rappresentabile graficamente.

Una volta implementata la libreria, si è passati ad una fase di test in cui
si sono valutate e confrontate qualità e performace con \textit{DEAP}.

Infine sono stati svolti dei test sulla qualità dei risultati prodotti. Qui
l'obbiettivo è stato quello di verificare che l'implementazione proposta fosse
in grado di produrre risultati comparabili con quelli ottenuti da \textit{DEAP}.

\clearpage

\chapter{Contesto}

Come anticipato nell'introduzione, il contesto in cui si va a collocare il
lavoro di tesi è molteplice. Abbiamo infatti a che fare con algoritmi genetici,
calcolo parallelo ed explainable AI.

\section{Algoritmi genetici}

Gli algoritmi genetici fanno parte delle possibili tecniche euristiche adottate
nell'ambito dell'ottimizzazione. Forniscono infatti soluzioni subottimali ma
in tempi generalmente inferiori rispetto ad algoritmi in grado di trovare
soluzioni ottime, ma con complessità più elevata.

In contesti in cui non è necessario trovare la soluzione ottima ma si cerca
solo di minimizzare (o massimizzare) una certa funzione obbiettivo, gli
algoritmi genetici sono una valida scelta.

Il comportamento di questi algoritmi tende a simulare l'evoluzione di una
popolazione di individui, i quali rappresentano possibili soluzioni, tramite
meccanismi basati sulla fitness. Ogni individuo è rappresentato da un cromosoma,
il quale ha generalmente la forma di una lista o di un vettore, i cui attributi
costituiscono la soluzione al problema o una codifica di essa. Nella loro forma
più comune sono composti da sei fasi:

\begin{enumerate}
	\item Generazione:
	\item Selezione:
	\item Crossover:
	\item Mutazione:
	\item Valutazione:
	\item Rimpiazzo:
\end{enumerate}


\clearpage

\chapter{Stato dell'arte}


\clearpage

\chapter{Metodologia}

\clearpage

\chapter{Sperimentazione}

\addcontentsline{toc}{chapter}{Conclusioni} % Capitolo non numerato
\chapter*{Conclusioni}


\bibliographystyle{plain}
\bibliography{bibliography} % Entries are in the refs.bib file
% \addbibresource{bibliography.bib}
% \printbibliography

\end{document}